\begin{claim}{C.13}{Proved under stated hypotheses}{Removing the small generators when filtered colimits are exact}\label{C13}
Let $\mathcal M$ satisfy the manuscript's (a)--(d), with the actual adjunction $\varphi\dashv\kappa$ assumed in (a). Suppose in addition that $\mathcal M$ has all small coproducts and exact filtered colimits (AB5). Then the canonical functor
\[
 \mathcal M[W^{-1}]\longrightarrow D_{\geq0}(k),\qquad X\longmapsto L\varphi(X),
\]
is an equivalence. We do not require a small projective generating category, projectivity of its $k$-Hom modules, local presentability, or an exact resolving functor on $\mathcal M$. This gives the underlying theorem in another precisely stated scope. It does not assert that AB5 follows from (a)--(d).
\end{claim}
\begin{proof}
We replace the missing exact resolving functor by a category of resolutions. Let $\mathcal R$ consist of nonnegative complexes $P$ of projective objects of $\mathcal M$ such that $H_iP=0$ for $i>0$. Give it weak equivalences
\[
 u\in W_{\mathcal R}\quad\Longleftrightarrow\quad
 \varphi(u)\text{ is a quasi-isomorphism},
\]
where $\varphi$ is applied termwise. A cofibration is a degreewise split monomorphism whose cokernel belongs to $\mathcal R$. We check the six axioms of an ABC cofibration category, in the convention of \citet[Definition 1.1.1]{RadulescuBanu2009}.

The zero object and closure under composition are immediate. Two-out-of-three holds by the definition of $W_{\mathcal R}$. The pushout of a cofibration exists degreewise. Its terms are projective, since after a degreewise splitting it is the direct sum of the other source and the cokernel. Its positive homology vanishes by the homology exact sequence. The resulting map is again a cofibration with the same cokernel. Termwise application of $\varphi$ preserves this split exact sequence, so acyclic cofibrations remain acyclic after pushout.

For factorization let $f:P\to Q$, $X=H_0P$, and $Y=H_0Q$. In $\mathcal M$ factor $H_0f$ as
\[
 X\xrightarrow{(H_0f,\,\alpha\otimes X)}
 Z=Y\oplus(J\otimes X)\xrightarrow{\mathrm{pr}_Y}Y.
\]
The first map is monic, and the second belongs to $W$ by (d). If $C$ is its cokernel, choose a projective resolution $R$ of $C$. The horseshoe construction produces a projective resolution $S$ of $Z$ and a degreewise split exact sequence $0\to P\to S\to R\to0$. We may simultaneously construct a chain map $g:S\to Q$ extending the given $f$ and lifting $Z\to Y$. Here are the relative lifting details: in degree zero choose the map on the projective summand $R_0$ through $Q_0\twoheadrightarrow Y$; after the maps in lower degrees have been chosen, the required next differential lands in the cycles of $Q$, and the projectivity of $R_n$ lifts it through the surjection from $Q_n$ onto those cycles. The already fixed map on $P_n$ is $f_n$. This induction gives $g$ with $g|_P=f$. Since $g$ induces $Z\to Y$ on $H_0$, it is in $W_{\mathcal R}$. Thus $P\to S\to Q$ is the required factorization.

Coproducts of resolutions are resolutions because AB5 implies AB4, and coproducts of projectives are projective. Coproducts of cofibrations remain degreewise split with resolution cokernel. Since $\varphi$ preserves coproducts and coproducts of $k$-modules are exact, coproducts of acyclic cofibrations are acyclic. Finally a countable sequence of cofibrations has a degreewise colimit whose terms are projective: after choosing successive splittings, each term is the direct sum of the initial projective and the successive projective cokernels. Exactness of filtered colimits shows that its positive homology and the positive homology of the relevant quotient vanish. Termwise $\varphi$ preserves this colimit, so a countable composite of acyclic cofibrations is acyclic. This verifies the remaining axioms. Every object of $\mathcal R$ is cofibrant.

The projection $H_0:\mathcal R\to\mathcal M$ induces an equivalence after inverting the maps inducing isomorphisms on $H_0$. Indeed every object has a projective resolution, every map lifts to a chain map between resolutions, and two lifts are chain homotopic. Chain homotopies become equalities on localization: their usual chain-cylinder has projective terms, positive homology zero, and projection inducing an $H_0$-isomorphism. Conversely a quasi-isomorphism between these bounded below projective resolutions is a chain-homotopy equivalence, by splitting its acyclic projective cone. The usual comparison argument therefore proves full faithfulness as well as essential surjectivity. After further localization this gives
\[
 \mathcal R[W_{\mathcal R}^{-1}]\simeq\mathcal M[W^{-1}]. \tag{C-res}\label{Cres}
\]
The same assertion holds diagramwise for every small indexing category $I$. To justify coherence, we use projective resolutions in $\mathcal M^I$. This diagram category has enough projectives: the left adjoints to evaluation send a projective $P$ to the diagram with value $\bigoplus_{I(i,j)}P$ at $j$, and sums of these give projective epimorphisms onto arbitrary diagrams. Their values are projective in $\mathcal M$. A projective resolution of the whole $I$-diagram maps to any given objectwise projective resolution of that diagram, and the same comparison and cylinder argument applies. This proves the diagramwise version of \eqref{Cres}, rather than merely an equivalence of the categories at the terminal index.

The functor
\[
 \Phi:\mathcal R\longrightarrow\Ch_{\geq0}(k),\qquad P\longmapsto\varphi(P),
\]
is now an exact functor of cofibration categories. It preserves the specified cofibrations, their pushouts, coproducts, and countable cofibration colimits, and takes weak equivalences to quasi-isomorphisms. Its terms are projective over $k$ because $\kappa$ is exact. ABC cofibration categories have homotopy colimits with their usual adjunction to restriction \citep[Theorem 9.6.3]{RadulescuBanu2009}. For an explicit preservation check, let $p_t:\Delta'I\to I$ be the direct replacement of a small indexing category used in \citet[Proposition 9.5.7 and the proof of Theorem 9.6.3]{RadulescuBanu2009}. Choose a restricted Reedy cofibrant resolution $Q\to p_t^*X$. Its colimit over $\Delta'I$ computes $\operatorname{hocolim}_I X$. The diagram $\Phi Q$ is again restricted, because $\Phi$ preserves weak equivalences, and Reedy cofibrant, because it preserves latching pushouts and cofibrations. The degree filtration of $\Delta'I$ builds its colimit from sums, pushouts, and a countable cofibration composite; all three are preserved by $\Phi$. Consequently $\Phi(\operatorname{colim}_{\Delta'I}Q)\cong\operatorname{colim}_{\Delta'I}\Phi Q$ supplies the required homotopy-colimit comparison. No right adjoint to $\Phi$, or unverified preservation of infinite constructions by a merely finite-exact functor, is used.

For completeness, the same homotopy-colimit structure is represented in $\mathcal M$ by all monomorphisms and $W$. The finite axioms follow from derived exact sequences and the displayed interval factorization. Coproducts preserve $L\varphi$, by taking sums of projective resolutions and using AB4. Filtered colimits also preserve $L\varphi$: for a filtered diagram choose a projective resolution in $\mathcal M^I$ as above. Its colimit is a projective resolution, since the colimit functor is exact and preserves projectives (its right adjoint, the constant diagram functor, is exact). Applying $\varphi$ commutes with that colimit, and homology commutes with filtered colimits in $\Mod k$. This proves the necessary weak-equivalence closure for the last two cofibration axioms. Thus $H_0$ preserves the indicated colimits and cofibrations, and the diagramwise equivalence transfers the homotopy-colimit comparison to $L\varphi$ on $\mathcal M[W^{-1}]$.

We can now repeat the realization argument with its comparison made explicit. If $K\in\Ch_{\geq0}(\mathcal M)$, resolve its Dold--Kan denormalization $\Gamma K$ projectively in the abelian category $s\mathcal M$. This gives a simplicial object of $\mathcal R$. Applying $\Phi$ and normalizing in the simplicial direction gives a first-quadrant double complex with terms $\varphi(Q_{p,q})$. Each $Q_{p,q}$ is projective in $\mathcal M$, and each vertical column resolves $K_p$: normalization is exact and takes projective objects of $s\mathcal M$ to projective objects of $\Ch_{\geq0}(\mathcal M)$. Additive functors commute with this normalization, because its normalized summand is the image of the standard split idempotent built from degeneracies. The total of $Q_{p,q}$ is consequently a projective resolution of $K$ in the derived sense. The finite-diagonal comparison gives
\[
 L\varphi\bigl(\operatorname{hocolim}\Gamma K\bigr)
 \simeq\Tot\bigl(\varphi(Q_{p,q})\bigr)
 \simeq L\varphi(K).
\]
As in Claim~\ref{C11}, the homotopy-colimit/constant-diagram counit for $\Delta^{\mathrm{op}}$ is invertible. The weak equivalences are saturated: they are closed under retracts by their definition, so the cofibration-category saturation criterion applies \citep[Theorem 7.2.7]{RadulescuBanu2009}. The induced functor $L\varphi$ on $\mathcal M[W^{-1}]$ is conservative: Brown's homotopy calculus of fractions represents every morphism by $X\to Z\leftarrow Y$ with the second map a weak equivalence; if its $L\varphi$-image is invertible, the first map also belongs to $W$ \citep[Theorem 6.4.4]{RadulescuBanu2009}. Consequently the simplicial enlargement localized at homotopy-colimit equivalences is precisely
\[
 D_{\geq0}(\mathcal M)[(L\varphi)^{-1}(\mathrm{Iso})^{-1}].
\]
Finally (c), coproduct preservation just proved, and retracts imply that $L\varphi\kappa P\to P$ is an isomorphism for every projective $k$-module $P$. Resolving a nonnegative complex by projective $k$-modules and using the first-quadrant comparison proves this counit invertible on all of $D_{\geq0}(k)$. The localized adjunction $L\varphi\dashv\kappa$ therefore gives the asserted canonical equivalence, by the unit--counit argument of Claim~\ref{C4}.
\end{proof}
